\documentclass[11pt,a4paper]{article}
\usepackage[T1]{fontenc}
\usepackage{lmodern}
\usepackage[margin=28mm]{geometry}
\usepackage{amsmath,amssymb,amsthm}
\usepackage{microtype}
\usepackage[colorlinks=true,linkcolor=blue,citecolor=blue,urlcolor=blue]{hyperref}
\newcommand{\Z}{\mathbb Z}
\newcommand{\R}{\mathbb R}
\newcommand{\E}{\mathbb E}
\newcommand{\indicator}{\mathbf 1}
\DeclareMathOperator{\Ext}{Ext}
\newtheorem{theorem}{Theorem}[section]
\newtheorem{lemma}[theorem]{Lemma}
\newtheorem{proposition}[theorem]{Proposition}
\numberwithin{equation}{section}
\title{Draws in the undirected trapping game on site-percolated lattices}
\author{Kaiwen Jin\\\small\href{mailto:yui.ui.siki@gmail.com}{yui.ui.siki@gmail.com}}
\date{13 October 2026}
\hypersetup{pdftitle={Draws in the undirected trapping game on site-percolated lattices},pdfauthor={Kaiwen Jin}}
\begin{document}
\maketitle
\begin{abstract}
Independently delete each vertex of the nearest-neighbor lattice
$\Z^d$ with probability $p$. Two players, knowing the entire resulting
board, move a token alternately to an undeleted neighbor that has never
been visited; a player with no legal move loses. A draw means that neither
player can force a win. We prove that the draw probability is zero in
dimension two for every $0<p<1$, whereas in every dimension $d\ge3$
it is positive for some strictly positive deletion probability.
The proof expresses the outcome through the effect of distant boundary
conditions on maximum independent sets. In the plane, a deterministic
boundary bound is incompatible with the fluctuations that a positive draw
probability would produce. In higher dimensions, contour shifts and a
uniform comparison of permuted deletion configurations retain the
influence of the two checkerboard boundary conditions. The comparison is
performed for all fixed geometries before a contour is selected from
the random board.
\end{abstract}

\section{Introduction}\label{sec:introduction}
The undirected trapping game, or Trap, is a game of perfect information
on a graph. A token starts at a specified vertex. The players alternate
moving it to a neighboring vertex that has not previously been visited,
and the player with no legal move loses. A finite board always has a
winner under optimal play. On an infinite board neither player need
have a strategy forcing a win; this outcome is called a draw.
The distinction is sensitive to the whole graph, since a strategy may
depend on arbitrarily distant parts of the board.

We consider the nearest-neighbor graph on $\Z^d$, with each vertex
independently deleted with probability $p\in(0,1)$. Both players know
the sampled graph before play begins. The starting vertex is the origin.
If it is deleted, the first player is declared the winner, so that this
event is not a draw. Let $D_d(p)$ denote the probability that neither
player can force a win. The question is whether the draw phase of the
undeleted lattice survives a positive density of deleted vertices.

Basu, Holroyd, Martin and W\"astlund \cite{BHMW} studied Trap on random
boards and related its finite outcomes to graph optimization.
For bipartite graphs their criterion can be expressed in terms of maximum
independent sets: the first mover loses at a vertex precisely when every
maximum independent set contains it. This relation suggests comparing
the two checkerboard boundary conditions of the lattice.
A persistent difference at the origin can then prevent either player
from forcing a win on the infinite board. The symmetric site-deletion
question is stated in the open problems of \cite{BHMW} and in AIM
Problem~574 \cite{AIM574}.
In the plane they also prove absence of draws in an asymmetric model:
for each fixed positive deletion probability on one parity, a sufficiently
small deletion probability on the other parity excludes draws
\cite[Theorem~1(ii), arXiv version~1, page~3]{BHMW}.
That result leaves equal positive deletion probabilities unresolved.

There is a related question for the hard-core model, which assigns
weight $\lambda^{|I|}$ to each independent set $I$ at activity
$\lambda$. Ayuso Ventura, Chiarini, Helmuth and Powell \cite{AVCHP}
prove uniqueness of its Gibbs measure on site-percolated $\Z^2$ at
every finite activity when the retention probability is less than one.
Their argument compares disorder fluctuations with the size of the
boundary and uses a spatial symmetry to interchange the checkerboards.
This result concerns finite activity. Trap instead depends on exact
maximum-cardinality sets, and passage to infinite activity alone does
not identify its winning strategies. A separate connection between
finite boundary conditions and the infinite game is therefore needed.

We establish the dimension-dependent picture for the original
undirected game:
\begin{theorem}\label{thm:main}
For every $p\in(0,1)$,
\[
 D_2(p)=0.
\]
For every integer $d\ge3$, there exists $p_d\in(0,1)$ such that
\[
 D_d(p_d)>0.
\]
\end{theorem}
The rules include self-avoidance of the token path and complete knowledge
of the board. The theorem makes no assertion about the location of a
transition or the exact set of deletion probabilities with draws.

The finite-to-infinite connection uses games in which reaching the
boundary permits an exit with a prescribed winner. These games have
an exact independent-set description. Their root influences are
monotone as the domain grows, and a draw is exactly a disagreement
between their limiting outcomes. In dimension two, we integrate these
influences into a finite-coordinate difference of independence numbers.
Its size is bounded by the boundary length. If draws had positive
probability, its martingale fluctuations would have a nonzero Gaussian
characteristic-function limit on the same scale. The deterministic
bound rules this out. This is a maximum-cardinality version of the
fluctuation-versus-boundary comparison in \cite{AVCHP}; we give the
limiting argument directly.

For $d\ge3$, we use contours surrounding regions with the wrong
checkerboard pattern. Shifting such regions by one lattice step and
filling the contour increases the number of occupied vertices.
To respect the deletions, the occupation shift is accompanied by a
finite permutation of the board coordinates and removal of the newly
occupied deleted sites. A contour chosen from the board gives a
board-dependent permutation, whose output need not be iid.
We instead compare all fixed admissible permutations simultaneously.
Dyadic majority approximations control their number at each scale,
and rare-deletion concentration controls their increments. The
resulting series converges precisely in dimensions above two.
The estimates make the two checkerboard boundary conditions force
opposite root occupations with positive probability uniformly in the box.
The monotone exit-game limits then give a draw on the infinite board.
Spatial shifts and disorder comparisons also appear in the
low-temperature framework of Chen, Zhou, Liu and Zhou \cite{CZLZ}.
Here the argument is given for exact independence numbers with the
original finite boundary restriction; no Gibbs-state theorem is used
to infer a game outcome.

Section~\ref{sec:games} establishes the finite boundary description
and its infinite-game limit. Section~\ref{sec:planar} proves the planar
assertion. Sections~\ref{sec:contours} and \ref{sec:random} give the
deterministic contour comparison and the uniform random estimate
needed in higher dimensions.

\section{Finite games and boundary conditions}\label{sec:games}
Write $\rho(v)=\sum_{i=1}^d v_i\pmod 2$ for the parity of a lattice vertex.
The player whose moves land on parity $\tau$ vertices is called the
$\tau$-player. At an even starting vertex the first mover is the $1$-player.
Let $\eta_v$ be one if $v$ is deleted and zero otherwise.
We use $X_v=1-\eta_v$ for the open indicator and $q=1-p$.

For a finite simple graph $G$, let $\alpha(G)$ be its independence number.
The finite-game criterion of Basu, Holroyd, Martin and W\"astlund
\cite[Proposition~5 and the following bipartite equivalence]{BHMW} gives
\begin{equation}\label{eq:finite-criterion}
 \text{the first mover loses from }v
 \quad\Longleftrightarrow\quad
 \alpha(G)-\alpha(G-v)=1
\end{equation}
when $G$ is bipartite and $v\in G$.
Indeed, their Vicious Trap criterion is that every maximum independent set
contains $v$. On a bipartite graph its additional deletion option removes
only vertices of the deleting player's destination parity, so its outcome
agrees with ordinary Trap. If a maximum independent set omits $v$, deleting
$v$ leaves $\alpha$ unchanged. Otherwise deleting $v$ lowers $\alpha$ by
exactly one. For a disconnected graph, apply the criterion to the component
of $v$ and use additivity of $\alpha$.
The statements used here are in the cited arXiv version~1,
pages~6--7; the Proposition~5 proof is on page~15.

For finite $V\subset\Z^d$, let
\[
 B(V)=\{v\in V:\text{ a lattice neighbor of }v\text{ is outside }V\},
 \qquad A_V^\tau=B(V)\cap\{\rho=1-\tau\}.
\]
In the exit-$\tau$ game a player standing at a boundary vertex may choose
an exit, ending the game with the $\tau$-player winning. This option is
offered without consulting the outside board. Set
\begin{equation}\label{eq:boundary-alpha}
 a_V^\tau(X)=\alpha\bigl(G[V\cap\{X=1\}\setminus A_V^\tau]\bigr),
 \qquad
 s_V^\tau(v;X)=a_V^\tau(X^{v=1})-a_V^\tau(X^{v=0}).
\end{equation}
Here $X^{v=i}$ changes only the coordinate at $v$.
Removing one vertex changes an independence number by at most one, so
$s_V^\tau\in\{0,1\}$.

\begin{lemma}\label{lem:exit}
At an open starting vertex $v\in V$, the exit-$\tau$ winner has parity
$\rho(v)$ exactly when $s_V^\tau(v;X)=1$.
\end{lemma}
\begin{proof}
At each $v\in B(V)$ of parity $1-\tau$, attach a private leaf of parity
$\tau$. At each boundary vertex of parity $\tau$, attach a private path
of two edges whose last vertex has parity $\tau$. Entering either gadget
has a forced continuation ending with a $\tau$-move, hence realizes the
declared exit winner. Retain all gadget vertices even when their attachment
vertex is closed.

A maximum independent set can omit every leaf attachment of parity
$1-\tau$: replace an occupied attachment by its private leaf, without
changing cardinality or creating an edge. Each leaf then contributes one.
A two-edge gadget contributes exactly one of its two private vertices,
whether or not its attachment is occupied. Thus the augmented graph has
independence number $|B(V)|+a_V^\tau(X)$.
This constant is unchanged when the starting vertex is removed, because
its gadgets remain. The augmented graph is finite, simple and bipartite;
\eqref{eq:finite-criterion} proves the assertion.
\end{proof}

\begin{lemma}\label{lem:infinite}
For every fixed board, at an even vertex the quantities $s_V^0(v)$
decrease and $s_V^1(v)$ increase as $V$ increases, with
$s_V^0(v)\ge s_V^1(v)$. At an odd vertex the inequalities reverse.
Their limits $s^\tau(v)$ exist and do not depend on the increasing
exhaustion. With $v$ forced open, put
\[
 h_v(X)=s^0(v;X)-s^1(v;X).
\]
Then $h_v$ ignores $X_v$ and equals $(-1)^{\rho(v)}$ on a draw at that
forced-open vertex, and zero otherwise. In particular,
\begin{equation}\label{eq:draw-influence}
 D_d(p)=q\,\E h_0.
\end{equation}
\end{lemma}
\begin{proof}
A strategy winning an exit game in which exits favor its opponent cannot
allow any exit along a compatible history. It uses only moves inside $V$,
and cannot reach a boundary vertex on the opponent's turn, since that
opponent could exit and win. Consequently it remains winning in every
larger domain and on the actual infinite board.

Conversely, a strategy forcing a finite win on the infinite board has a
uniform finite bound on its compatible histories. Its history tree is
finitely branching; if depths were unbounded, successively choosing a
child with arbitrarily deep descendants would produce an infinite
compatible play. A domain containing this finite history tree and a
one-site buffer realizes the win even with exits awarded to the opponent.

These statements imply monotonicity in the domain. Awarding exits to a player can only improve that player's outcome, which
gives the comparison of the two boundary conditions. Every finite domain
eventually lies in any increasing exhaustion, so the limits agree across
exhaustions. At an even vertex the pair of limits can only be $(0,0)$,
$(1,1)$ or $(1,0)$: the first two give an actual finite win by the odd or
even player, respectively, by the correspondence just proved. The last
gives neither win, hence a draw. The odd-vertex case reverses the roles.
Forcing $v$ open makes $h_v$ independent of its own coordinate. An actually
closed origin is not a draw, giving \eqref{eq:draw-influence}.
\end{proof}

The indicators above are measurable: finite-game outcomes are functions
of finitely many coordinates, and their limits are pointwise limits.
For $T_vX$ defined by $(T_vX)_w=X_{w+v}$, the exhaustion independence also
gives
\begin{equation}\label{eq:translation-h}
 h_v(X)=(-1)^{\rho(v)}h_0(T_vX).
\end{equation}
The iid board law is invariant under translations and under
$\theta(v)=(1-v_1,v_2,\ldots,v_d)$, which interchanges the two parities.

We will use elementary conditional averages on this countable product
space. A few facts behind the limiting averages are useful to record.
Finite-coordinate conditioning is a finite sum on coordinate atoms.
For a bounded $f$, its averages on nested finite coordinate sets are
orthogonal projections in $L^2$; their squared norms increase and remain
at most $\E f^2$. The Pythagorean identity makes them Cauchy.

A bounded $L^2$-Cauchy sequence has a bounded $L^2$ limit without any
additional probabilistic theorem: select a subsequence whose consecutive
$L^2$ distances are at most $2^{-2j}$. Markov's bound makes the probability
of a difference exceeding $2^{-j}$ at most $2^{-2j}$. A tail union bound
shows that only finitely many such exceptions occur almost surely; the
subsequence is then pointwise Cauchy. Bounded convergence and the original
Cauchy property give convergence of the whole sequence.

Bounded finite-coordinate functions are dense in $L^2$ of the
coordinate-generated sigma algebra. For indicators, approximable events
contain the cylinder algebra and are closed under complements and finite
unions, by clipping approximants to $[0,1]$ and taking maxima. Countable
unions follow by first using a finite partial union and continuity of
probability. Quantization gives the bounded-function statement. The same
argument holds for any coordinate subfamily. The projection limit
therefore satisfies the conditional averaging identity for that subfamily.
These facts also give convergence of averages on growing cubes to the
original bounded function.

All uses of bounded convergence below can be seen directly from
$\E|U_n-U|\le\epsilon+M\Pr(|U_n-U|>\epsilon)$ for uniformly bounded
differences; countable additivity gives convergence in probability from
pointwise convergence. Symmetry of the iid law follows first on cylinder
events and then on their generated sigma algebra.


\section{Absence of planar draws}\label{sec:planar}
In this section $d=2$ and
$V_n=[-n+1,n]^2\cap\Z^2$. These squares are preserved by $\theta$.
The infinite-game influence $h_v$ will be detected by finite-coordinate
fluctuations of independence numbers.

For finite $L$ and $n$ large enough that $L\subset V_n$, define
\[
 \Psi_{n,L}^\tau(X)=a_{V_n}^\tau(X)-a_{V_n}^\tau(X^{L=1}).
\]
Forcing sites open changes cardinality by at most one per site, so
$|\Psi_{n,L}^\tau|\le |L|$. Telescoping over the finitely many changed
coordinates expresses this difference as a sum of signed $s_{V_n}^\tau(v)$
in fixed finitely perturbed boards. Lemma~\ref{lem:infinite} makes each
binary summand eventually constant. Consequently $\Psi_{n,L}^\tau$
has a pointwise limit $\Psi_L^\tau$ on every board. Put
\[
 H_{n,L}=\Psi_{n,L}^0-\Psi_{n,L}^1,\qquad
 H_L=\Psi_L^0-\Psi_L^1,\qquad F_L=\E[H_L\mid X_L].
\]
These are bounded measurable functions, and $F_L$ depends only on $X_L$.
Since the all-open reference is unchanged by toggling $v\in L$,
bounded convergence and finite-coordinate conditioning give
\begin{equation}\label{eq:toggle-F}
 F_L(X^{v=1})-F_L(X^{v=0})
 =\E[h_v\mid X_{L\setminus\{v\}}].
\end{equation}

\begin{lemma}\label{lem:surface}
For $L=V_j$,
\begin{equation}\label{eq:surface}
 \E F_L=0,\qquad |F_L|\le 8j+4.
\end{equation}
\end{lemma}
\begin{proof}
Remove the annulus $C_j=V_{j+1}\setminus V_j$ from the already reduced
graphs in $V_n$, with $n\ge j+2$. Its cardinality is $8j+4$. The removal
separates the unrestricted inner graph on $L$ from the outer graph,
which retains its original boundary restriction at $B(V_n)$.
No new boundary restriction is placed along the annulus.

Conditional on $X_L$, the inner independence number is identical for the
two parities. Reflection $\theta$ interchanges the outer boundary
restrictions and preserves the law of the outside coordinates. The
conditional expected difference of the two independence numbers after
removal is therefore zero. Removal decreases each number by an amount
in $[0,|C_j|]$, so their difference changes by at most $|C_j|$.
It follows that
\[
 \left|\E[a_{V_n}^0(X)-a_{V_n}^1(X)\mid X_L]\right|\le |C_j|.
\]
For $X^{L=1}$, the expected parity difference is zero by reflection even
conditional on $X_L$: that reference ignores $X_L$, and its all-open
inner square is reflection invariant. Subtract the reference and pass
to the limit using the bound $|H_{n,L}|\le2|L|$. This proves the pointwise
bound on $F_L$. Reflecting the entire board changes $H_L$ to $-H_L$,
which proves $\E F_L=0$.
\end{proof}

Order the lattice lexicographically and enumerate $L=V_j$ as
$v_1,\ldots,v_N$, $N=4j^2$. Write
$\mathcal F_i=\sigma(X_{v_1},\ldots,X_{v_i})$.
Equation~\eqref{eq:toggle-F} gives the Doob differences
\begin{equation}\label{eq:doob}
 F_L=\sum_{i=1}^N Y_{i,j},\qquad
 Y_{i,j}=(X_{v_i}-q)a_{i,j},\qquad
 a_{i,j}=\E[h_{v_i}\mid\mathcal F_{i-1}].
\end{equation}
To verify the coefficient, the two possible conditional means after
exposing $X_{v_i}$ differ by the average of \eqref{eq:toggle-F} over
unexposed sites. The field $h_{v_i}$ ignores its own coordinate.
Thus $|Y_{i,j}|\le1$ and
\[
 \E[Y_{i,j}^2\mid\mathcal F_{i-1}]=q(1-q)a_{i,j}^2.
\]

Let $\mathcal P_v=\sigma(X_w:w\prec v)$ be the full lexicographic past and
put $a_v=\E[h_v\mid\mathcal P_v]$. The bounded projection construction in
Section~\ref{sec:games} defines this average. Its uniqueness and
\eqref{eq:translation-h} allow versions
$a_v=(-1)^{\rho(v)}a_0\circ T_v$. In particular $a_v^2$ is stationary.
Set
\[
 b^2=q(1-q)\E a_0^2.
\]

\begin{lemma}\label{lem:variance}
As $j\to\infty$,
\begin{equation}\label{eq:variance-limit}
 Q_j:=\frac1N\sum_{i=1}^N
 \E[Y_{i,j}^2\mid\mathcal F_{i-1}]
 \longrightarrow b^2\quad\text{in probability},
 \qquad b^2\ge q(1-q)(\E h_0)^2.
\end{equation}
\end{lemma}
\begin{proof}
Let $a_v^{(r)}$ condition on past sites in $v+[-r,r]^2$.
The projection construction gives
$\|a_v^{(r)}-a_v\|_2\to0$, uniformly in $v$ after its parity sign.
If the radius-$r$ cube lies in $L$, the finite revealed past contains
these coordinates and lies in $\mathcal P_v$. Nested projection errors
give
$\|a_{i,j}-a_{v_i}\|_2\le\|a_{v_i}^{(r)}-a_{v_i}\|_2$.
The exceptional sites form a boundary fraction tending to zero at fixed
$r$. Boundedness and $|x^2-y^2|\le2|x-y|$ for $|x|,|y|\le1$ show
\[
 \frac1N\sum_i |a_{i,j}^2-a_{v_i}^2|\longrightarrow0
 \quad\text{in }L^1
\]
by first sending $j$, then $r$, to infinity.

For the spatial average of $a_v^2$, approximate it by
$g_v^{(r)}=\E[a_v^2\mid X_{v+[-r,r]^2}]$.
These bounded stationary variables are independent when their cubes
are disjoint. The variance of their average is at most $(4r+1)^2/N$.
The average approximation error has expectation no greater than
$\E|a_0^2-g_0^{(r)}|$, which tends to zero by finite-coordinate density.
Thus $N^{-1}\sum_{v\in L}a_v^2\to\E a_0^2$ in $L^1$.
This proves the limit for $Q_j$.
Finally $\E a_0=\E h_0$, and nonnegativity of
$\E(a_0-\E a_0)^2$ gives the lower bound.
\end{proof}

For fixed $t\in\R$, put $Z_i=Y_{i,j}/\sqrt N$ and
$r_i=\E[Z_i^2\mid\mathcal F_{i-1}]$.
The power-series remainder for the exponential, conditional mean zero,
and $|Z_i|\le N^{-1/2}$ give
\[
 \E[\exp(itZ_i+t^2r_i/2)\mid\mathcal F_{i-1}]
 =1+O_t(N^{-3/2}).
\]
Here $r_i\le N^{-1}$ and the remainder is uniform in $i,j$.
For
\[
 T_k=\exp\left(it\sum_{i\le k}Z_i+
                    \frac{t^2}{2}\sum_{i\le k}r_i\right)
\]
we have $|T_k|\le e^{t^2/2}$. Conditional averaging and telescoping
therefore give $\E T_N=1+O_t(N^{-1/2})$.
By Lemma~\ref{lem:variance}, the bounded factor $e^{t^2Q_j/2}$
converges in $L^1$ to $e^{t^2b^2/2}$. Replacing that factor in the
expectation, with an error bounded by its $L^1$ difference, yields
\begin{equation}\label{eq:characteristic}
 \E e^{itF_{V_j}/\sqrt{|V_j|}}\longrightarrow e^{-t^2b^2/2}.
\end{equation}
This argument uses predictable compensators throughout.

\begin{proposition}\label{prop:planar}
For every $p\in(0,1)$, $D_2(p)=0$.
\end{proposition}
\begin{proof}
Lemma~\ref{lem:surface} bounds $U_j=F_{V_j}/\sqrt{|V_j|}$ in $[-6,6]$.
Let $\phi_j(t)=\E e^{itU_j}$.
For any positive integer $m$ and any real $s\ne0$, a binomial expansion gives
\[
 \frac{(-1)^m}{s^{2m}}\sum_{k=0}^{2m}(-1)^k
       \binom{2m}{k}\phi_j((m-k)s)
 =\E\left(\frac{2\sin(sU_j/2)}s\right)^{2m}
 \le6^{2m}.
\]
At fixed $m,s$, use \eqref{eq:characteristic} in this finite sum.
Then let $s\to0$ in the resulting analytic expression.
The binomial difference annihilates degrees below $2m$; the leading
coefficient gives $(2m-1)!!\,b^{2m}\le6^{2m}$ for every $m$.
If $b>0$ this is impossible: at least $\lfloor m/2\rfloor$ factors in
$(2m-1)!!$ are at least $m$, so its $2m$-th root tends to infinity.
Hence $b=0$. Since $q(1-q)>0$, \eqref{eq:variance-limit} gives
$\E h_0=0$, and \eqref{eq:draw-influence} proves the result.
\end{proof}


\section{Contours and comparison of maximum sets}\label{sec:contours}
We next work in arbitrary dimension $d\ge2$.
For $\tau\in\{0,1\}$ put $b^\tau(v)=\indicator_{\{\rho(v)=\tau\}}$.
Let $C_z=2z+\{0,1\}^d$, $z\in\Z^d$, be the coarse cells, and write
$\widehat S=\bigcup_{z\in S}C_z$.
All coarse neighborhoods use integer shifts.
A cell box $\Omega=\widehat\Lambda$ has $\Lambda$ a finite rectangular
coarse box.

An occupation $n:\Z^d\to\{0,1\}$ is admissible if it is independent on
the full nearest-neighbor lattice, equals $b^\tau$ outside $\Omega$,
and is zero at deleted sites inside $\Omega$.
Its interior set $I=\{v\in\Omega:n_v=1\}$ is exactly an independent set
of the reduced graph in \eqref{eq:boundary-alpha}.
Indeed, an occupied outside neighbor forbids precisely the
opposite-parity internal boundary sites.
The exterior checkerboard is this finite boundary reference; deletion
constraints are imposed only in the free box.
Write $\alpha_\Omega^\tau(\eta)=a_\Omega^\tau(1-\eta)$ for its maximum
interior cardinality.

For such an occupation let $\mathcal D$ be the edges with two empty
endpoints. This set is finite, since the outside occupation is a
checkerboard. Let $B$ be the set of coarse cells at sup distance at most
one from a cell containing an endpoint of $\mathcal D$, and put
$\Gamma=B+[-2,2]^d$. Its components use star adjacency, meaning sup
distance one between distinct coarse cells.

For finite coarse $S$, let $\Ext(S)$ be the unique infinite
nearest-neighbor component of $S^c$, and let
$H(S)=\Z^d\setminus\Ext(S)$ be its filled hull.
Uniqueness follows because the exterior of a large box is connected
in $d\ge2$, and every infinite complementary component reaches that
exterior. The other components are finite.

We use the following lattice consequence of Tim\'ar's boundary theorem
\cite[Theorem~3]{Timar}: for a finite star-connected $S$, the
nearest-neighbor boundary adjacent to $S$ in each complementary
nearest-neighbor component $R$ is star connected.
The cited version~2 states the exterior case and the case inside a cubic
box. A bounded global $R$ fits, with $S$, into a translated larger cube;
translation preserves both adjacencies and the boundary, giving the
stated consequence.

\begin{lemma}\label{lem:collar}
Let $S$ be a star component of $\Gamma$. For every nearest-neighbor
component $R$ of $S^c$, the cells
\[
 D_R=\{z\in R:z\text{ has a nearest neighbor in }S\}
\]
and the adjacent cells of $S$ have one common checkerboard label
$\tau_R$. Thus the fine occupation there agrees with $b^{\tau_R}$.
\end{lemma}
\begin{proof}
A cell outside $B$, with its radius-one cell neighborhood, contains no
empty-edge endpoint. Independence excludes occupied pairs, so on this
connected fine box every edge has exactly one occupied endpoint.
Its occupation is one checkerboard. Neighborhoods of star-adjacent good
cells overlap, so their labels agree.

Every cell in $D_R$ is outside $\Gamma$, since otherwise it would join
the component $S$. Each adjacent cell of $S$ is outside $B$: a cell of
$B$ has its radius-two neighborhood in $\Gamma$, hence in its own
component. Tim\'ar's result connects $D_R$ by star steps. The good-cell
overlaps identify all its labels and those of the adjacent $S$ cells.
Defects farther inside $R$ play no role.
\end{proof}

\begin{lemma}\label{lem:outer-selection}
If $n_v\ne b^\tau(v)$ for some $v\in\Omega$, there is a component $S$
of $\Gamma$ with exterior label $\tau$ and with the coarse cell of $v$
in $H(S)$.
\end{lemma}
\begin{proof}
A nearest-neighbor component of $\Gamma^c$ has one checkerboard label.
Its infinite component has label $\tau$, because it reaches the fixed
outside region.
If the root cell $z$ is outside $\Gamma$, its label is wrong, hence
its good component $R_0$ is finite. The exterior nearest-neighbor
boundary of $R_0$, visible from infinity in $R_0^c$, lies in $\Gamma$
and is star connected by Tim\'ar's result. It lies in one component $S$.
Every infinite simple path from $R_0$ has a last departure from this
finite set; the next vertex is on that visible boundary. Hence $S$
separates $z$ from infinity and $z\in H(S)$.
If $z\in\Gamma$, begin with its own component.

If $S$ has the wrong exterior label, choose an exterior collar cell
$x\in\Ext(S)$. It is outside $\Gamma$ and its good component $R$ has
the wrong label, so $R$ is finite. Its visible exterior boundary lies
in a component $S'$ of $\Gamma$ and separates $x$ from infinity.
Thus $x$ is in a bounded component of $(S')^c$. The component $S'$ cannot
be $S$, because $x\in\Ext(S)$.

Distinct $\Gamma$ components are at sup distance at least two.
A star step within $S$ can be interpolated by a nearest-neighbor path
in its unit cube avoiding $S'$. Together with the edge from $x$ to $S$,
these paths place all of $S$ in the bounded component of $(S')^c$
containing $x$. Conversely $R\subset\Ext(S)$, and a point of its visible
boundary connects to $R$ while avoiding $S$. Interpolating star paths
in $S'$ while avoiding $S$ puts all of $S'$ in $\Ext(S)$.

It follows that $\Ext(S')\subset\Ext(S)$ and
$H(S)\subsetneq H(S')$; strictness uses
$S'\subset H(S')\setminus H(S)$.
Repeating selects distinct components of the finite $\Gamma$, so the
ascent terminates with the required exterior label and preserves enclosure
of the root.
\end{proof}

\begin{lemma}\label{lem:transport}
Let $S$ have exterior label $\tau$ and cardinality $k$.
Let $A$ be the fine union of bounded complementary components whose
collar label is $1-\tau$, and let $u=e_1$ be a fine unit vector.
There is a finite permutation $\pi_A$, supported inside the original
$\Omega$, such that for $\eta'_v=\eta_{\pi_A(v)}$,
\begin{equation}\label{eq:projected-comparison}
 \alpha_\Omega^\tau(\eta')\ge |I|+
 c_d k-\sum_{v\in T(S)\cap\Omega}\eta_v,\qquad
 c_d=\frac1{4d\,15^d},\quad T(S)=\widehat{S+[-1,1]^d}.
\end{equation}
This assertion holds for every admissible $I$, with no probabilistic
hypothesis.
\end{lemma}
\begin{proof}
Define
\[
 n'_v=\begin{cases}
 b^\tau(v),&v\in\widehat S,\\
 n_{v-u},&v\in A,\\
 n_v,&\text{otherwise}.
 \end{cases}
\]
On each line parallel to $u$, write the maximal fine intervals of $A$
as $[a,b]$. Call $a-u$ the head and $b$ the tail of such an interval.

Every opposite-label bounded component is inside the free cell box.
Otherwise an outward coarse ray from a cell outside that box, avoiding
the box, first meets $S$ after passing through fixed $\tau$ sites.
Its last cell is on the component collar and would give label $\tau$,
a contradiction. The interval heads one fine step outside $A$ are
inside $\Omega$ too: their cells are adjacent $S$ collar cells with
label $1-\tau$, incompatible with the fixed exterior checkerboard.
Thus $n'$ keeps the prescribed outside values.

An empty-edge endpoint cell outside $S$ belongs to $B$. Its radius-two
neighborhood lies in another $\Gamma$ component, so its sup distance
from $S$ is at least four. Translating an empty edge by one fine step
keeps it in the same complementary component.
Near a shifted component boundary the old label $1-\tau$ becomes
$\tau$, since $b^{1-\tau}(v-u)=b^\tau(v)$.
The unchanged collars already have label $\tau$.
Hence $n'$ is independent, no crossing empty edge is created, and
empty edges outside $\widehat S$ are unchanged or translated
bijectively inside their own components.

Let $m$ count old empty edges in $\widehat S$.
No old empty edge crosses its boundary, by the radius-two buffer around
its endpoint cells in $B$.
All these $m$ edges disappear. For independent occupations the empty-edge
indicator is $1-n_x-n_y$. Summing its difference over the finitely many
affected edges gives
\begin{equation}\label{eq:edge-identity}
 m=2d\sum_{v\in\Omega}(n'_v-n_v).
\end{equation}
Every fine vertex has degree $2d$, and outside changes vanish.
Moreover $k\le5^d|B\cap S|$. Each contributing empty edge supplies at most
$2\cdot3^d$ cells of $B$, and those source edges lie wholly in $\widehat S$.
Thus $k\le2\,15^d m$ and $m/(2d)\ge c_d k$.

For these intervals, set
\[
 \pi_A(v)=v-u\quad(v\in[a,b]),\qquad
 \pi_A(a-u)=b,
\]
and fix all other sites. The extended intervals $[a-u,b]$ are disjoint
cycles and contained in $\Omega$, so $\pi_A$ is a bijection of its board
coordinates. If $v\in A$ and $n'_v=1$, its source $\pi_A(v)=v-u$ was
occupied and open. Off $A\cup\widehat S$ both transformations are the
identity. Therefore
\[
 C=\{v\in\Omega:n'_v=1,\ \eta'_v=1\}\subset\widehat S.
\]
The distinct sources $\pi_A(C)$ are sites of $S$ or interval tails
adjacent to $S$, and lie in $T(S)\cap\Omega$. Consequently
$|C|\le\sum_{T(S)\cap\Omega}\eta_v$.

Set $n^*_v=0$ on $C$ and $n^*_v=n'_v$ elsewhere.
This occupation is independent, open under $\eta'$ inside $\Omega$,
and has the original outside boundary. Its interior cardinality is
$|I|+m/(2d)-|C|$ by \eqref{eq:edge-identity}.
It is an allowed set for $\alpha_\Omega^\tau(\eta')$, proving
\eqref{eq:projected-comparison}.
\end{proof}

For a maximum original $I$, \eqref{eq:projected-comparison} compares
two maximum-size problems on jointly transformed boards. The map used
by an actual contour depends on the board. The next section bounds all
fixed admissible geometries simultaneously before making that selection.


\section{Uniform comparison and higher-dimensional draws}
\label{sec:random}
Fix the cell box $\Omega$, its boundary parity $\tau$, a coarse root $z_0$,
and an integer $k\ge1$ before sampling the board.
Let $\mathcal A_{\Omega,z_0,k}$ be the distinct fine sets $A$ obtained
by choosing any bounded complementary components of a star-connected
coarse $S$ with $|S|=k$ and $z_0\in H(S)$, subject to
$A$ and all its interval heads lying inside $\Omega$.
The class is deterministic and finite. It contains the empty set.
All geometries from Lemmas~\ref{lem:outer-selection} and
\ref{lem:transport} belong to it.

Define the process
\[
 Z_A(\eta)=\alpha_\Omega^\tau(P_A\eta)-\alpha_\Omega^\tau(\eta),
 \qquad (P_A\eta)_v=\eta_{\pi_A(v)}.
\]
For fixed $A$, $P_A$ is a permutation of iid coordinates inside $\Omega$.
We need a bound holding for the entire fixed class, since the geometry
of an actual maximum set is selected from the board.

\subsection{Geometry and majority approximations}
Every coordinate span of a star-connected $S$ with $k$ cells is at most
$k-1$. Its filled hull lies in its coordinate bounding box: outside that
box an outward ray avoids $S$. Thus all its cells are at sup distance
at most $k-1$ from $z_0$.
Every bounded complementary component has a nearest-neighbor edge to
$S$, so there are at most $2dk$ such components.
With $\Delta=3^d-1$, rooted depth-first walks give at most
\begin{equation}\label{eq:support-count}
 (2k-1)^d\Delta^{2(k-1)}
\end{equation}
supports enclosing $z_0$.
Every fine cut edge of $A$ meets $\widehat S$, since a neighbor outside
$S$ would be in the same selected complementary component.
Consequently its edge boundary and interval-head count are at most
$B_d k$, where $B_d=d2^d$.

Tile the fine lattice by aligned cubes of side $l=2^\ell$ and let
$U_\ell(A)$ be the union of cubes at least half filled by $A$.
Ties are included.

\begin{lemma}\label{lem:entropy}
There are dimension-only constants $F_d,E_d>0$ such that, uniformly
over $\Omega,z_0,k,\ell$,
\begin{align}
 |A\triangle U_\ell(A)|&\le B_dkl,\label{eq:majority-error}\\
 \#\{\text{cube-lattice cut edges of }U_\ell(A)\}
 &\le \left\lfloor F_d k/l^{d-1}\right\rfloor,\label{eq:coarse-cuts}\\
 \log\#\{U_\ell(A):A\in\mathcal A_{\Omega,z_0,k}\}
 &\le E_d(1+\ell)k/l^{d-1}.\label{eq:entropy}
\end{align}
One may take $F_d=8dB_d$.
\end{lemma}
\begin{proof}
Join ordered pairs in a rectangular lattice block by changing coordinates
in the order $1,\ldots,d$. For an edge in direction $i$, path congestion
is at most $L_i$ times the block volume, where $L_i$ is its side length.
In a cube of $N=l^d$ sites let $r$ be the smaller of the two label counts.
The ordered-pair difference sum is $2r(N-r)\ge rN$, and congestion is
at most $lN$. Hence $r\le l$ times its internal cut-edge count.
Summing disjoint cubes proves \eqref{eq:majority-error}.

Two adjacent cubes with different majority labels have at least $N/2$
sites of each type in their union. The difference sum is at least $N^2$,
while congestion is at most $4l^{d+1}$. There are therefore at least
$l^{d-1}/4$ original cut edges in this rectangle.
Each fine edge belongs to at most $2d$ such rectangles, giving
\eqref{eq:coarse-cuts} with $F_d=8dB_d$.

For completeness we give the position count, including disconnected and
nested interfaces. Put $M=\lfloor F_d k/l^{d-1}\rfloor$.
Address a cube-lattice cut edge by $(z,i)$ joining the cubes at $lz$ and
$l(z+e_i)$. Join addresses when their base points are at sup distance
at most three. This fixed graph has degree at most $W_d-1$, with
$W_d=d7^d$. Split the cut set into components of sizes
$q_1,\ldots,q_m$, total $Q\le M$, and choose one root edge in each.

The rectangle argument locates a fine cut edge near every root.
That cut edge meets $\widehat S$, so the integer fine root address $lz$
is within $\ell^1$ distance $(2d+1)l$ of some $2s$, $s\in S$.
Order roots by first visits to these attachment points along a
depth-first traversal of a spanning tree of $S$.
The fine traversal length is at most $4d(k-1)$, and its start is within
$2d(k-1)$ of $2z_0$. Thus total root-list jump length is at most
\[
 6dk+2(2d+1)lm\le K_d k,\qquad
 K_d=6d+2(2d+1)F_d,
\]
because $m\le M$ and $d\ge2$.
The support traversal proves existence of this ordering; it need not
be encoded.

A code records the ordered composition of $Q$, integer root increments
with total absolute sum at most $K=K_dk$, root orientations, and rooted
component depth-first walks. For given $Q,m$ their number is at most
\[
 \binom{Q-1}{m-1}\,W_d^{2Q}\,d^m\,
 2^{dm}\binom{K+dm}{dm}.
\]
The walks recover the cut set. A finite cube union is determined by its
cut edges and label zero at infinity: the modulo-two difference of any
two such labelings is constant on the connected cube lattice.
Therefore valid codes decode uniquely; every attained majority union
has a code. Invalid codes only enlarge this upper bound.

The inequality $\binom{n}{r}\le(en/r)^r$ follows from the $r$-th
term of $e^r=\sum_j r^j/j!$.
The function $m\log(1+K/m)$ is increasing.
Summing compositions over $Q\le M$ gives, when $M\ge1$,
\[
 \#\{U_\ell(A)\}\le
 1+\exp\{a_d M+dM\log(1+K/M)\},
 \quad a_d=\log2+\log d+2\log W_d+d(1+\log2).
\]
Since $\lfloor x\rfloor\ge x/2$ for $x\ge1$,
$K/M\le (2K_d/F_d)l^{d-1}$.
This proves \eqref{eq:entropy}; for example one can choose
\[
 E_d=F_d\left[\log2+a_d+d\log(1+2K_d/F_d)
                     +d(d-1)\log2\right].
\]
If $M=0$, a nonempty finite cube union would have a cut edge, so every
majority union is empty and its count is one.
All counts used only geometric properties shared by every admissible $A$;
the original-box restriction takes a subclass.
\end{proof}

\subsection{Centered increments for fixed permutations}
Changing one deletion coordinate changes $\alpha_\Omega^\tau$ by at
most one. For fixed $A,A'$ put
\[
 D=\{v:\pi_A(v)\ne\pi_{A'}(v)\},\qquad
 R=\pi_A(D)=\pi_{A'}(D).
\]
The equality follows from bijectivity and agreement on $D^c$.
Conditional on the inputs in $\Omega\setminus R$, the common output
coordinates outside $D$ are fixed, while the coordinates in $D$ are
two permutations of the same unexposed iid inputs. Thus
\begin{equation}\label{eq:conditional-center}
 \E[Z_A-Z_{A'}\mid\eta_{\Omega\setminus R}]=0.
\end{equation}
The boundary graph is the same function of each output array; no graph
permutation invariance is required.
Changing one input in $R$ changes this difference by at most two.

For $0<p<1/4$ let
\[
 \nu(p)=\max\{\sqrt p,4/\log(1/p)\}.
\]
If $U=\eta-p$, then $\E U=0$, $\E U^2\le p$, $|U|\le1$.
For $|t|\le h=\tfrac12\log(1/p)$, the exponential-series remainder gives
$\E e^{tU}\le1+\tfrac12t^2e^{|t|}p\le e^{t^2\sqrt p/2}$.
For $|t|>h$, it gives $\E e^{tU}\le e^{|t|}\le e^{t^2/h}$.
Hence $\E e^{tU}\le e^{\nu(p)t^2/2}$ for every real $t$.
Exposing the $|R|$ inputs in \eqref{eq:conditional-center} yields
increments $(\eta_i-p)c_i$ with past-measurable $|c_i|\le2$.
Iterating the conditional exponential bound and optimizing Markov's
inequality therefore gives
\begin{equation}\label{eq:increment-tail}
 \Pr(|Z_A-Z_{A'}|\ge t)
 \le2\exp\left\{-\frac{t^2}{8\nu(p)|D|}\right\}.
\end{equation}
For $D=\varnothing$ the difference is zero.
The maps agree as shifts on $A\cap A'$ and as identities away from
their symmetric difference and interval heads, so
\begin{equation}\label{eq:permutation-metric}
 |D|\le |A\triangle A'|+\#\text{heads}(A)+\#\text{heads}(A')
 \le |A\triangle A'|+2B_dk.
\end{equation}

\subsection{A finite chain of actual admissible representatives}
For each majority class select an actual $A$ in the class, before
sampling $\eta$, using increasing cardinality and then a fixed
lexicographic order. Denote the representative of the class of $A$ by
$r_\ell(A)$. The empty set is selected for an empty majority union.
At $\ell=0$ the representative is $A$ itself. Let $L_k$ be the first
integer with $2^{L_k(d-1)}>F_dk$. By Lemma~\ref{lem:entropy},
$r_{L_k}(A)=\varnothing$ and $Z_{r_{L_k}(A)}=0$.

The approximation bound gives
$|r_\ell(A)\triangle A|\le2B_dk2^\ell$, hence
$|r_\ell(A)\triangle r_{\ell+1}(A)|\le6B_dk2^\ell$.
Equations~\eqref{eq:increment-tail}--\eqref{eq:permutation-metric}
bound each consecutive increment by
\[
 2\exp\{-t^2/(V_d\nu(p)k2^\ell)\},\qquad V_d=64B_d.
\]
The log number of attained pairs is at most the sum of the two cover
entropies, bounded by
$H_d(1+\ell)k/2^{\ell(d-1)}$, $H_d=2E_d$.
No consistency of parent majority classes is needed.

Now assume $d\ge3$. Put
\[
 A_d=\sqrt{2V_dH_d},\qquad
 \beta=(\sqrt2-1)/8,\qquad
 S_d=\sum_{\ell\ge0}\sqrt{1+\ell}\,2^{-\ell(d-2)/2}<\infty.
\]
Choose $p$ so small that $A_d\sqrt{\nu(p)}S_d\le c_d/8$.
Use thresholds, for $0\le\ell<L_k$,
\[
 t_\ell=A_d\sqrt{\nu(p)}\,k\sqrt{1+\ell}\,2^{-\ell(d-2)/2}
             +\beta c_dk\,2^{(\ell-L_k)/2}.
\]
Their sum is at most $c_dk/4$. The square of the first term pays twice
the log pair count in the exponential bound; the square of the second
leaves $\beta^2c_d^2k/(V_d\nu(p)2^{L_k})$.
A union bound over pairs and levels gives
\begin{equation}\label{eq:uniform-tail}
 \begin{aligned}
 \Pr\left(\max_{A\in\mathcal A_{\Omega,z_0,k}}Z_A\ge c_dk/2\right)
 &\le2L_k\exp\left\{-\frac{\beta^2c_d^2k}
                         {V_d\nu(p)2^{L_k}}\right\}\\
 &\le C_d\log(2k)\exp\{-b_dk^{(d-2)/(d-1)}/\nu(p)\}.
 \end{aligned}
\end{equation}
with $C_d,b_d>0$ depending only on $d$.
Indeed, outside the increment failures the chain telescopes to the empty
representative and bounds every $|Z_A|$ by $c_dk/4$.
Also $2^{L_k}\le2(F_dk)^{1/(d-1)}$ and
$L_k\le1+\log_2(F_dk)/(d-1)$.
All representatives stay in the original admissible class; their majority
cube unions themselves are never used as board permutations.

The sum of the right-hand side of \eqref{eq:uniform-tail} over $k\ge1$
tends to zero as $p\downarrow0$.
For a fixed small $\nu_0$, its dyadic block $2^j\le k<2^{j+1}$ is
bounded by $C(j+1)2^j e^{-b2^{j(d-2)/(d-1)}/\nu_0}$, a summable series.
Smaller $\nu$ have this common tail bound, while each fixed term tends
to zero. A finite head and uniform tail justify the limit.

\subsection{Forcing the root and passing to the infinite game}
For any support $S$ of size $k$,
$|T(S)|\le6^dk$. Put $a=c_d/2$ and $M_d=6^d$.
For $p<a/(eM_d)$, exponential Markov with
$e^t=a/(M_dp)$ gives
\[
 \Pr\left(\sum_{T(S)\cap\Omega}\eta_v\ge ak\right)
 \le(eM_dp/a)^{ak}.
\]
Multiplying by \eqref{eq:support-count} and summing over $k$ yields a
quantity tending to zero with $p$. For small enough $p$ the exponential
ratio $\Delta^2(eM_dp/a)^a$ is below $e^{-2}$, which gives a common
polynomial-times-exponential summable bound.

If a maximum set has a wrong root, Lemma~\ref{lem:outer-selection}
selects an enclosing exterior-$\tau$ support, and
\eqref{eq:projected-comparison} gives
$Z_A\ge c_dk-\sum_{T(S)\cap\Omega}\eta_v$ for its admissible $A$.
Thus either $Z_A\ge c_dk/2$ or the last sum is at least $c_dk/2$.
The preceding bounds hold for all fixed supports and admissible $A$
simultaneously. They can therefore be used for this selected geometry
without asserting an iid law for its adaptively transformed board.

For any $\delta>0$ choose a strictly positive sufficiently small $p$
so that the total exception probability is below $\delta$, uniformly
in $\Omega$ and $\tau$. Every even-boundary maximum then occupies the
even origin with probability at least $1-\delta$, and every odd-boundary
maximum omits it with probability at least $1-\delta$.
The first event forces an open origin and $s_\Omega^0(0)=1$.
The second implies $s_\Omega^1(0)=0$ on that open-origin event.
It follows that
\[
 \Pr(X_0=1,\ s_\Omega^0(0)=1,\ s_\Omega^1(0)=0)\ge1-2\delta.
\]
Along the exhaustion $\Omega_n=[-2n,2n+1]^d$ these weaker
cardinality-increment events decrease by Lemma~\ref{lem:infinite}.
Their intersection is exactly an infinite-board draw at an open origin.
Continuity of probability from above gives $D_d(p)\ge1-2\delta$.
Choose $\delta=1/4$ to obtain a positive draw probability for each
fixed $d\ge3$.

\begin{proof}[Proof of Theorem~\ref{thm:main}]
Proposition~\ref{prop:planar} proves the first assertion for every
$p\in(0,1)$. The preceding uniform comparison and infinite-game passage
prove the second for each integer $d\ge3$ with a strictly positive
dimension-dependent deletion probability. Both use the original
undirected rules and the stipulated exclusion of a deleted origin
from draws.
\end{proof}

\section*{AI-use statement}
Codex was used to develop candidate arguments, perform finite diagnostic
checks, draft and revise the manuscript, and check proof details and
citations. The proof presented here is analytic; the diagnostic computations
are not premises of the theorem. Independent mathematical review is requested.

\begin{thebibliography}{9}
\bibitem{BHMW}
R.~Basu, A.~E.~Holroyd, J.~B.~Martin and J.~W\"astlund,
\emph{Trapping games on random boards},
Annals of Applied Probability \textbf{26} (2016), no.~6, 3727--3753.
\href{https://doi.org/10.1214/16-AAP1190}{doi:10.1214/16-AAP1190}.
Statement and proposition numbering used here:
\href{https://arxiv.org/abs/1505.07485v1}{arXiv:1505.07485v1} (2015).

\bibitem{Timar}
\'A.~Tim\'ar,
\emph{Boundary-connectivity via graph theory},
Proceedings of the American Mathematical Society \textbf{141}
(2013), 475--480.
\href{https://doi.org/10.1090/S0002-9939-2012-11333-4}
{doi:10.1090/S0002-9939-2012-11333-4}.
Version used here:
\href{https://arxiv.org/abs/0711.1713v2}{arXiv:0711.1713v2} (2011).

\bibitem{AVCHP}
I.~Ayuso Ventura, L.~Chiarini, T.~Helmuth and E.~Powell,
\emph{Imry-Ma phenomenon for the hard-core model on $\Z^2$},
\href{https://arxiv.org/abs/2601.05798v1}{arXiv:2601.05798v1} (2026).

\bibitem{CZLZ}
Y.~Chen, J.~Zhou, R.~Liu and H.-J.~Zhou,
\emph{The stability of long-range order in disordered systems:
A generalized Ding-Zhuang argument},
\href{https://arxiv.org/abs/2507.11445v1}{arXiv:2507.11445v1} (2025).

\bibitem{AIM574}
AIM Open Problems,
\emph{574. Draws in the undirected trapping game on percolation clusters},
problem statement at the pinned revision.
\href{https://github.com/MColbrook/AIM/blob/c9929805cece044705abaaebe1cebddac1ba11a6/problems/574-trap-draws-percolation.md}
{Pinned AIM statement}.
\end{thebibliography}
\end{document}
